\section{Introduction}

Public-sector endpoint fleets --- the workstations, servers, domain controllers, and
special-purpose hosts operated by government agencies --- accumulate disclosed
vulnerabilities faster than constrained operations teams can remediate them. The
publication rate of Common Vulnerabilities and Exposures (CVEs) continues to grow,
while each maintenance window admits only a bounded number of changes under
change-advisory, blackout, and approval constraints. The practical question is
therefore not \textit{which vulnerabilities are severe} but *which vulnerability-host
pairs should consume the next unit of limited remediation capacity\textit{, and }whether
that decision can be explained and reviewed afterward*.

Existing signals each address part of the problem. CVSS \cite{first_cvss_v31}
quantifies intrinsic severity but is a weak predictor of real-world exploitation
\cite{allodi2014comparing}. EPSS \cite{jacobs2021epss} estimates exploitation likelihood
but is asset-agnostic. The CISA KEV catalog \cite{cisa_kev} and Binding Operational
Directive 22-01 \cite{cisa_bod2201} impose deadlines on confirmed-exploited CVEs but
cover a small fraction of the backlog. Commercial risk-based vulnerability
management (RBVM) products combine such signals with proprietary asset context, but
their scoring is closed and not independently reproducible. Prior machine-learning
prioritization research advances modeling but typically reports a single tuned
result on private or non-reproducible data and rarely produces per-decision audit
evidence.

We do not claim a better predictor. We ask whether an open, reproducible,
audit-evidence-producing benchmark can be built that ranks vulnerability-host pairs
under public-sector-shaped operational constraints, and whether combining exploit
intelligence with asset criticality and endpoint exposure improves prioritization
over EPSS-style baselines --- treated as an open empirical question. Our
contributions are (i) an integrated framework (feed normalization, exploit
enrichment, synthetic endpoint telemetry, asset-criticality and local-exposure
modeling, vulnerability-host pair construction, scoring/ranking strategies, a
capacity-constrained scheduling simulation, and append-only hash-chained audit
records); (ii) a deterministic, public-sector-shaped synthetic evaluation with a
frozen 30-seed artifact and an explicit freeze/verify integrity protocol; and (iii)
an honest empirical baseline. We state plainly that, under the current toy fixtures
and uncalibrated weights, the proposed model is statistically indistinguishable
from the EPSS baseline and does not outperform a random ordering. The value of the
work is the reproducible, falsifiable evaluation structure, not a superiority claim.
